\section{ARIMA, diagnostics et validation temporelle}\label{sec:arima}

\subsection{Identification sur le log-prix}



Un ARIMA(p,d,q) combine une différenciation d’ordre d, p retards autorégressifs et q retards d’innovations. Avec d=1, il décrit la variation du log-prix plutôt que son niveau persistant. Le drift représente une variation moyenne par séance, sans imposer une hausse quotidienne. La représentation ci-dessous utilise le polynôme autorégressif \(\Phi\), le polynôme de moyenne mobile \(\Theta\) et l’opérateur de retard B.



\begin{equation}\label{eq:arima} \Phi(B)(1-B)^d y_t=c+\Theta(B)\varepsilon_t.\end{equation}



Le choix d=1 vient des diagnostics TRAIN et est confirmé par trois vérifications \sourcepath{ndiffs()} en R. L’ACF/PACF ne présente pas de coupure suffisamment nette pour imposer un seul ordre. Cinq candidats manuels couvrent le random walk sans drift, le random walk avec drift et les modèles (1,1,0), (0,1,1), (1,1,1) avec drift. Ils permettent d’évaluer progressivement la dépendance moyenne et la complexité.



Les recherches automatiques sont bornées à p,q\(\leq\)3, avec d fixé, sans saisonnalité et sans approximation de vraisemblance. \sourcepath{auto.arima} recherche un modèle selon les critères et contraintes fournis \cite{R6}. En R, la proposition automatique est ARIMA(2,1,2) avec drift. En Python, \sourcepath{auto_arima} propose (3,1,3) sans drift ; une règle de parcimonie et de diagnostics retient provisoirement (1,1,1) avec drift, proche en AICc et plus simple. Les propositions ne sont pas assimilées à des gagnants prédictifs.



Les implémentations consomment le même TRAIN. Une sensibilité du (2,1,2) Python initialisé par les coefficients R rapproche sa vraisemblance de celle de R. Cela documente un effet d’initialisation et d’optima locaux ; cette vérification dépend de R et ne devient ni une réplication indépendante ni une preuve de maximum global. Les tentatives rejetées et avertissements de convergence restent dans les sorties.



\subsection{Adéquation de la moyenne et limites des innovations}



Le R (2,1,2) retenu provisoirement donne des p-values Ljung-Box de 0,869791 et 0,784385 aux retards 10 et 20. Le Python (1,1,1) donne 0,067465 et 0,118122. Ces non-rejets concernent l’autocorrélation linéaire testée, sans établir des innovations indépendantes. Jarque-Bera rejette fortement la normalité ; les p-values numériques nulles correspondent à un sous-flux. ARCH-LM rejette aux retards 5, 10 et 20 dans les deux langages. La moyenne est donc mieux représentée que la variance et les queues.



Le gate SARIMA reste fermé : la STL exploratoire, les retards saisonniers et les vérifications R ne fournissent pas l’évidence conjointe exigée. Aucun SARIMA n’est estimé. La figure R de diagnostics figure en annexe A ; les conclusions formelles proviennent des CSV et de leur convention de résidus.



Sources : \sourcepath{outputs/tables/phase3_arima_diagnostic_summary.csv}, \sourcepath{outputs/tables/phase3_seasonality_gate.csv}, \sourcepath{docs/phase3a_arima_identification.md}.





% Documentary display only; edit preparation definitions for a refresh.
% Source CSVs: outputs/tables/phase3_arima_diagnostic_summary.csv
\begin{table}[!htbp]
\centering
\caption{Diagnostics des modèles de moyenne provisoires de Phase 3A. TRAIN ; p-values Ljung-Box et ARCH-LM, sans unité. Six décimales ou six chiffres significatifs. Jarque-Bera rejette dans les deux langages ; zéro numérique par sous-flux.}
\label{tab:arima_diagnostics}
\small
\setlength{\tabcolsep}{4pt}
\renewcommand{\arraystretch}{1.16}
\begin{tabularx}{\linewidth}{lXrrr}
\toprule
\textbf{Langue} & \textbf{Modèle avec drift} & \textbf{LB(10)} & \textbf{LB(20)} & \textbf{ARCH-LM(5)} \\
\midrule
R & ARIMA(2,1,2) & 0,869791 & 0,784385 & \(1{,}50217\times 10^{-12}\) \\
Python & ARIMA(1,1,1) & 0,067465 & 0,118122 & \(1{,}09996\times 10^{-13}\) \\
\bottomrule
\end{tabularx}
\par\smallskip
\begin{minipage}{\linewidth}\footnotesize Sources : outputs/tables/phase3\_arima\_diagnostic\_summary.csv.
\end{minipage}
\end{table}


\subsection{Validation rolling-origin exclusivement sur TRAIN}



La Phase 3B déplace la question de l’ajustement vers la prévision. À chaque origine, les modèles sont estimés sur un préfixe croissant, puis évalués à cinq et dix séances. La fenêtre initiale contient 1386 lignes ; l’origine avance de dix séances. Les 59 origines vont du 2023-07-06 au 2025-10-27. Toutes les cibles restent dans TRAIN. La liste des candidats ayant été construite sur TRAIN complet, il s’agit d’une validation interne de sélection, sans prétendre à une évaluation entièrement imbriquée du processus d’identification.



\historicalfigure{phase3b_rmse_by_model_h10.png}{RMSE de validation ARIMA à dix séances}{Comparaison interne R/Python à h=10. Source : \sourcepath{outputs/figures/phase3b/phase3b_rmse_by_model_h10.png}. Périmètre : 59 origines TRAIN planifiées, 2023-07-06--2025-10-27 ; certains fits Python incomplets. Unité : USD ajustés. Les scores de panels incomplets ne sont pas directement comparables aux scores complets.}{0.32}



La figure~\ref{fig:phase3b_rmse_by_model_h10.png} montre des gains possibles à dix séances. Ils doivent être confrontés aux deux horizons, aux erreurs absolues et à la couverture. Python ne conserve que 53 fits utilisables pour (1,1,1) et 54 pour (2,1,2), contre 59 pour Naive ; R couvre les 59. Les comparaisons de gain sont calculées sur les origines appariées, et les fits non convergents ne sont pas remplacés silencieusement.



La règle définie avant les résultats exige une convergence complète et au moins 1 \% de gain en RMSE ET MAE aux deux horizons dans les deux langages. Aucun challenger ne satisfait toutes les conditions. La spécification verrouillée est donc Naive, ARIMA(0,1,0) sans drift, qui prévoit exactement le dernier prix observé. Le warning Ljung-Box du random walk reste conservé : le choix prédictif ne transforme pas ses résidus en bruit blanc validé.



Enfin, le protocole final inclut explicitement ARIMA(2,1,2) avec drift comme challenger pédagogique issu de la Phase 3A. Cette inclusion ne remplace pas le lock Naive de Phase 3B. Elle permet une comparaison finale pré-déclarée, sans rechercher un nouvel ordre sur le holdout. Qualité d’ajustement, choix interne et recommandation finale restent trois décisions distinctes.



Sources : \sourcepath{outputs/tables/phase3b_arima_validation_metrics.csv}, \sourcepath{outputs/tables/phase3b_baseline_comparison.csv}, \sourcepath{outputs/tables/phase3b_locked_arima_model.csv}, \sourcepath{docs/phase7_final_evaluation_protocol.md}.
% Documentary display only; edit preparation definitions for a refresh.
% Source CSVs: outputs/tables/phase3b_arima_validation_metrics.csv
\begin{table}[!htbp]
\centering
\caption{Validation interne ARIMA : sous-ensemble pédagogique Naive/(1,1,1)/(2,1,2). h en séances ; RMSE/MAE en USD ajustés, trois décimales. Scores sur origines disponibles : panels Python incomplets, non appariés à Naive dans ce tableau ; les gains de la règle utilisent les panels appariés.}
\label{tab:arima_validation}
\small
\setlength{\tabcolsep}{4pt}
\renewcommand{\arraystretch}{1.16}
\begin{tabularx}{\linewidth}{lXrrrr}
\toprule
\textbf{Langue} & \textbf{Candidat} & \textbf{h} & \textbf{n / prévu} & \textbf{RMSE} & \textbf{MAE} \\
\midrule
Python & (1,1,1) + drift & 5 & 53 / 59 & 7,838 & 6,032 \\
Python & (1,1,1) + drift & 10 & 53 / 59 & 8,780 & 6,496 \\
Python & (2,1,2) + drift & 5 & 54 / 59 & 7,351 & 5,434 \\
Python & (2,1,2) + drift & 10 & 54 / 59 & 8,672 & 6,619 \\
Python & Naive & 5 & 59 / 59 & 7,724 & 5,856 \\
Python & Naive & 10 & 59 / 59 & 9,369 & 7,253 \\
R & (1,1,1) + drift & 5 & 59 / 59 & 7,848 & 5,934 \\
R & (1,1,1) + drift & 10 & 59 / 59 & 9,081 & 6,731 \\
R & (2,1,2) + drift & 5 & 59 / 59 & 7,847 & 5,938 \\
R & (2,1,2) + drift & 10 & 59 / 59 & 9,080 & 6,729 \\
R & Naive & 5 & 59 / 59 & 7,724 & 5,856 \\
R & Naive & 10 & 59 / 59 & 9,369 & 7,253 \\
\bottomrule
\end{tabularx}
\par\smallskip
\begin{minipage}{\linewidth}\footnotesize Sources : \sourcepath{outputs/tables/phase3b_arima_validation_metrics.csv}.
\end{minipage}
\end{table}


Les couvertures et scores ci-dessus sont repris des CSV de validation interne, sans nouveau fit. La règle de remplacement utilise les erreurs appariées et conserve les échecs de convergence ; ce tableau de synthèse ne redéfinit pas cette règle.
